\section{Discussion}\label{sec:discussion}
\subsection{Matrix-free rSVD}
\subsubsection{Asymptotic analysis}
To explain the computational times and asymptotic behavior in Sections~\ref{sec:computational_time} and \ref{sec:computational_breakdown}, we examine the cost of forming $K$ and of Algorithm~\ref{alg:discrete_rsvd} in more detail.

For a 2D elliptic problem, the Cholesky factorization of $A$ and $M_\partial$ can be done in $\mathcal{O}(N^{3/2})$ and $\mathcal{O}(N_b^{3/2})$ computations, respectively (Section~\ref{sec:initial_matrix_factorizations}). The resulting triangular systems can be solved at the respective costs $\mathcal{O}(N \log N)$ and $\mathcal{O}(N_b \log N_b)$. Since $N_b \ll N$, we may disregard the factorization and triangular solves using $M_\partial$. 

The cost of computing $K = T A^{-1} M$ involves solving $N_b$ linear systems $A v_i = T^T e_i$ (Section~\ref{sec:approximating}), each system costing $\mathcal{O}(N \log N)$ using a sparse Cholesky factorization of $A$. With $N_b \sim \sqrt{N}$, the cost therefore scales as $\mathcal{O}(N^{3/2} \log N)$, explaining the scaling of $t_K$ observed in the left panel of Figure~\ref{fig:rsvd_computational_time_comp}. The SVD of dense matrix $K \in \mathbb{R}^{N_b \times N}$ costs $\mathcal{O}(N_b^2 N) = \mathcal{O}(N^2)$ \cite{golubMatrixComputations4th2013}, which is larger than $\mathcal{O}(N^{3/2} \log N)$, explaining $t_{\text{SVD}}$ in Figure~\ref{fig:rsvd_computational_time_comp}.

For the matrix-free rSVD (Algorithm~\ref{alg:discrete_rsvd}), factorizing $A$ costs $\mathcal{O}(N^{3/2})$, and is done once. Each application of $K$ and $K^*$ costs $\mathcal{O}(N \log N)$, for a total of $\mathcal{O}(l N \log N)$. The SVD of the $l \times N$ matrix $B$ costs $\mathcal{O}(l^2 N)$ \cite{golubMatrixComputations4th2013}, and the cost of the QR factorization of the small $N_b \times l$ matrix $Y$ is $\mathcal{O}(l^2 N_b)$ \cite{golubMatrixComputations4th2013}, and may be disregarded. Hence, assuming two dimensions and an efficient sparse Cholesky factorization of $A$, Algorithm~\ref{alg:discrete_rsvd} costs
\begin{equation}\label{eq:rsvd_cost}
    \underbrace{\mathcal{O}(N^{3/2})}_{\text{setup}}
    + \underbrace{\mathcal{O}(l N \log N)}_{\substack{\text{forward and} \\ \text{adjoint solves}}}
    + \underbrace{\mathcal{O}(l^2 N)}_{\text{SVD of } B}.
\end{equation}

This is consistent with the results in the right panel of Figure~\ref{fig:rsvd_computational_time_comp}: $t_{\text{rSVD}}$ scales near-linearly on the observed range, matching the near-linear scaling of $\mathcal{O}(N \log N)$. Even if $\mathcal{O}(N^{3/2})$ dominates $\mathcal{O}(N \log N)$ at large $N$, the hidden constant of the big-O notation makes the scaling be more similar to $\mathcal{O}(N \log N)$ (see Figure~\ref{fig:rsvd_computational_breakdown}), explaining the near-linear scaling in Figure~\ref{fig:rsvd_computational_time_comp}.

Equation~\eqref{eq:rsvd_cost} also explains the cost breakdown in Figure~\ref{fig:rsvd_computational_breakdown}. For larger $k$, a larger share is taken up by the SVD stage, as predicted by $\mathcal{O}(l^2 N)$. For $k = 200$ ($l = k + p$), the largest considered target rank, we observed that as $N$ increases, the forward and adjoint solves start taking up a larger share of the time compared to the SVD of $B$, also consistent with Equation~\eqref{eq:rsvd_cost}. The cost $\mathcal{O}(l^2 N_b)$ of the QR-factorization $QR = Y$ is negligible, matching the results for all $k$ and $N$.


\subsubsection{Overhead of the adjoint formulation}
The adjoint formulation of Algorithm~\ref{alg:discrete_rsvd_Kstar} has some computational overhead compared to Algorithm~\ref{alg:discrete_rsvd}. The setup stage requires an additional Cholesky factorization of the $N \times N$ matrix $M$, which costs $\mathcal{O}(N^{3/2})$. For the QR factorization, $Y$ is now $N \times l$, so this costs $\mathcal{O}(l^2 N)$, compared to $\mathcal{O}(l^2 \sqrt{N})$ for rSVD on $K$. Computing $B = Q^T K^*$ requires solving $\tilde q_i M = q_i$ $l$ times, costing $\mathcal{O}(l N \log N)$ (compared to $\mathcal{O}(l \sqrt{N} \log \sqrt{N})$, which was negligible). The SVD of $B \in \mathbb{R}^{l \times N_b}$ costs $\mathcal{O}(l^2 \sqrt{N})$, which is cheaper than $\mathcal{O}(l^2 N)$ for rSVD on $K$. Overall, the asymptotic cost of Algorithm~\ref{alg:discrete_rsvd_Kstar} is the same as Algorithm~\ref{alg:discrete_rsvd}, which is given in Equation~\eqref{eq:rsvd_cost}.

In Table~\ref{tab:computational_time_adjoint} ($t_1/t$), we see approximately 30\% computational overhead for $k = 10$, 20\% for $k = 50$, and 10\% for $k = 200$. This matches the theoretical overhead, which reduces for larger $k$ when the cost $\mathcal{O}(l^2 N)$ of the SVD in Algorithm~\ref{alg:discrete_rsvd} takes up a larger share.

To recover the SVD of $K$ from the SVD of $K^*$ (Section~\ref{sec:reconvering_the_svd}), QR factorizations of $A \in \mathbb{R}^{N_b \times k}$ and $B \in \mathbb{R}^{N \times k}$ costs $\mathcal{O}(k^2 \sqrt{N}) + \mathcal{O}(k^2 N) = \mathcal{O}(k^2 N)$ \cite{golubMatrixComputations4th2013}. The SVD of $R_A \bar \Sigma R_b \in \mathbb{R}^{k \times k}$ is $\mathcal{O}(k^3)$ \cite{golubMatrixComputations4th2013}, which is negligible compared to $\mathcal{O}(k^2 N)$, since $k \ll N$. This is in correspondence with the reported computational overhead $t_2 / t$ in Table~\ref{tab:computational_time_adjoint}: as $k$ increases, the extra cost $\mathcal{O}(k^2 N)$ of the recovery procedure and the reduced cost $\mathcal{O}(l^2 \sqrt{N})$ of the SVD of $B$ balance out, leading to approximately 30\% extra computational time using rSVD on $K^*$ and recovery of the SVD, compared to the direct approach of rSVD on $K$.


\subsubsection{Analysis of reconstruction error}\label{sec:analysis_of_reconstruction_error}
At target rank $k = 50$, rSVD reconstruction errors were within 0.5\% to 4\% of the exact Tikhonov solution in Euclidean norm and AUC-IoU (4\% to 39\% in EMD), with visually indistinguishable reconstructions (Section~\ref{sec:target_rank_reconstruction}, Figure~\ref{fig:test_problems_selected}).

To understand these results, we inspect the structure of $Y = K \Psi$. Following the setup of Halko et al. \cite[Sec.~9]{halkoFindingStructureRandomness2011}, let $K = U \Sigma V^T$ be the SVD of $K$, and partition $\Sigma$ as $\Sigma = \operatorname{diag}(\Sigma_1,\, \Sigma_2)$, where $\Sigma_1$ contains the $k$ first singular values of $K$, and $\Sigma_2$ the remainder. Correspondingly, partition the right singular values as $V = [V_1\,\, V_2]$, and write $K$ as
\begin{equation}\label{eq:rsvd_partitioned_1}
    K = U \begin{bmatrix}
    \Sigma_1 & \\ & \Sigma_2
    \end{bmatrix}
    \begin{bmatrix}
    V_1^T \\ V_2^T
    \end{bmatrix}.
\end{equation}
Right-multiplying by $\Psi$ and setting $\Psi_1 = V_1^T \Psi$ and $\Psi_2 = V_2^T \Psi$ gives
\begin{equation*}\label{eq:rsvd_partitioned_2}
    K = U \begin{bmatrix}
    \Sigma_1 \Psi_1 \\
    \Sigma_2 \Psi_2
    \end{bmatrix}.
\end{equation*}
Intuitively, $\Sigma_1 \Psi_1$ represents the actions of $K$, while $\Sigma_2 \Psi_2$ represents the perturbation. By Theorem~10.5 in \cite{halkoFindingStructureRandomness2011}, if $\Psi \in \mathbb{R}^{N \times l}$ is a random Gaussian matrix, then
\begin{equation}\label{eq:rsvd_error_bound}
    \mathbb{E}\| K - K_k \|_F \le
    \gamma_k \sqrt{\sum_{i = k+1}^{N_b} \sigma_i^2},
\end{equation}
where $\mathbb{E}$ denotes the expected value, $\gamma_k = \sqrt{1 + k / (p - 1)}$, and $K_k$ is the rank-$k$ approximation obtained from Algorithm~\ref{alg:discrete_rsvd}.

When the singular values of $K$ decay rapidly, the perturbation $\Sigma_2 \Psi_2$ remains low and the sum in Equation~\eqref{eq:rsvd_error_bound} is small, resulting in a good rank-$k$ approximation. If the singular values decay more slowly, $\sum_{i=k+1}^{N_b} \sigma_i^2$ will be larger, and the expected error increases.

Since $K$ has rapidly decaying singular values
(Figure~\ref{fig:singular_values}), a target rank
$k \ll \operatorname{rank}(K) = N_b$ yields an accurate rank-$k$
approximation: the main actions of $K$ are captured in $\Sigma_1 \Psi_1$. This argument transfers to the weighted Tikhonov problem in Equation~\eqref{eq:discrete_inverse_problem}. Both $K$ and $W$ are approximated accurately at $k = 50$, keeping the reconstruction error within 4\% (Euclidean, AUC-IoU). Since the random sampling introduces some noise, the EMD shows a higher error, as it is more sensitive to low-amplitude noise.


\subsubsection{Reconstruction error using the adjoint operator}
\begin{figure*}[thbp]
    \centerline{\includegraphics[width=1.0\textwidth]{true_singular_values.png}}
    \caption{Comparison of the normalized singular values $\bar \sigma_i(K)$ and $\bar \sigma_i(K^*)$ of the forward and adjoint operators, respectively. Left: Problem~I and II. Right: Problem~III.
    }
    \label{fig:true_singular_values}
\end{figure*}
Using rSVD on $K$ gives slightly lower reconstruction error than rSVD on $K^*$ (Section~\ref{sec:reconstruction_error_using_the_adjoint_operator}). The quality of the random sampling $Y = K\Psi$ versus $Y' = K^* \Psi'$ depends on the spectral decay of the operators. Figure~\ref{fig:true_singular_values} shows that the singular values of $K^*$ have a slower decay than those of $K$, and we therefore expect a higher approximation error by Equation~\eqref{eq:rsvd_error_bound}. The difference originates in the mass matrix factors of the relation $K^* = M^{-1} K^T M_\partial$ (Equation~\eqref{eq:K_star_def}).


\subsubsection{Effect of the test vector distribution}
\begin{figure*}[htbp]
    \centerline{\includegraphics[width=1.0\textwidth]{singular_vectors_of_K.png}}
    \caption{The first four right singular vectors $v_1,\dots,v_4$ (showed in the same order) of the forward operator $K$ from Problem~I.}
    \label{fig:singular_vectors_of_K}
\end{figure*}
Some test vector distributions slightly reduce reconstruction error, with the ``Peak'' test vectors giving lower EMD than Gaussian or Rademacher (Section~\ref{sec:test_vector_distribution_effect_on_error}).

Consider a single sample $y = K \psi = U \Sigma V^T \psi$. In the best case, $\psi = a v_1$ with $a \neq 0$: then $V^T(a v_1) = a e_1$, $K(a v_1) = a \sigma_1 u_1$, and the QR factorization gives $Q = u_1$. More generally, choosing $k$ linearly independent $\psi_i \in \operatorname{span}\{ v_1,\dots,v_k \}$ gives $\operatorname{span}(Q) = \operatorname{span}\{u_1,\dots,u_k\}$ and the rank-$k$ projection $QQ^T K$ recovers the exact rank-$k$ SVD of $K$.

Figure~\ref{fig:singular_vectors_of_K} plots the first four right singular vectors of $K$. We see that they are smooth, more closely matching the ``Peak'' test vectors than the Gaussian or Rade\-macher test vectors, explaining the small improvement in error.


\subsubsection{Independence from mesh resolution}
The reconstruction error is independent of the mesh resolution $N$ for fixed $k$ (Section~\ref{sec:mesh_resolution_effect_on_reconstruction_error}). The PDE solution at the boundary is smooth, so finer discretization produces more samples of the same smooth function without introducing new information about the source.

The continuous forward operator $\mathcal{K}$ has rapidly decaying singular values (Section~\ref{sec:ill_posedness}), and the discrete $K$ inherits this decay. As $N$ increases and $\operatorname{rank}(K) = N_b$ grows, the addition singular values are small with little new information.

The target rank $k$ can therefore be selected on a coarse mesh, where rSVD is inexpensive to run, and reused unchanged as the mesh is refined. This is a practical method for tuning $k$ when the final computation on a fine mesh is expensive.


\subsubsection{Estimation of singular values}
Let $\hat P = QQ^T$ with $Q$ produced by Algorithm~\ref{alg:discrete_rsvd}. This is an orthogonal projection onto $\operatorname{span}(Q)$, and $\hat P K$ is the projected matrix from which the rSVD approximation is obtained. The projector satisfies $\| \hat P z \| \le \| z \|$, meaning it shrinks vectors, so $\| \hat P K x  \| \le \|K x \|$ for every $x$. Consequently, $\hat P K$ stretches every vector at most as much as $K$ does, and singular values measure exactly this stretching. This is what we saw inSection~\ref{sec:approximation_of_singular_values}. By choosing a higher oversampling parameter $p$, $\operatorname{span}(Q)$ more closely matches $\operatorname{span}(K)$ (in particular, the $k$ first left singular vectors), and the underestimation reduces. Algorithm~\ref{alg:discrete_rsvd_Kstar} gave a worse approximation (Section~\ref{sec:reconstruction_error_using_the_adjoint_operator}), so the underestimation increases.

The captured subspace $\operatorname{span}(Q)$ depends on the random test vectors $\psi_i$, and so does the degree of underestimation. This is an active area of research. Lemma~4.2 and Theorem~4.3 of Gu \cite{guSubspaceIterationRandomization2015} provide lower bounds for the power-iteration method, which reduces to the plain rSVD used in Algorithm~\ref{alg:discrete_rsvd} when $q = 0$.


\subsubsection{Reconstructions under noisy data}
Choosing $k$ too large worsens reconstruction quality on noisy data (Section~\ref{sec:accuracy_on_noisy_data}), and the target rank acts as a regularization parameter in the same manner as the truncation rank of tSVD.

For a noisy measurement $\tilde y = y + \epsilon$, $y = K x$, we can write the minimum-norm least-squares solution as $x^* = K^+ \tilde y = K^+ y + K^+ \epsilon$ (Section~\ref{sec:the_pseudoinverse}). Expanding in terms of the SVD components,
\begin{equation*}
    x^* = \sum_{i=1}^{N_b}
    \left(\frac{u_i^T y + u_i^T \epsilon}{\sigma_i}\right) v_i,
\end{equation*}
we see that the $i$'th component depends on the exact term $u_i^T y$ and the noise term $u_i^T \epsilon$. Components where $|u_i^T y| \le |u_i^T \epsilon|$ are dominated by noise, and the $\tfrac{1}{\sigma_i}$ factor amplifies this contribution. By the Picard condition \cite{hansen2010discrete}, $u_i^T y$ decays to zero as $i$ increases while $u_i^T \epsilon$ stays approximately constant, so for large $i$ every component is noise-dominated and amplified by a growing $\tfrac{1}{\sigma_i}$.

By truncating the $N_b - k$ terms of the solution, we remove the noisy SVD components. We wish to choose $k$ such that $|u_k^T y| > |u_k^T \epsilon|$ while $|u_{k+1}^T y| < |u_{k+1}^T \epsilon|$, removing only the noise-dominated SVD components $k+1,\, k+2,\dots,\, N_b$.

For the weighted Tikhonov solution, the reconstruction already has some regularization through $\lambda$, but the argument transfers. We found that the optimal $k$ values are approximately equal for both the rSVD and the tSVD. The rank is problem-dependent, and for our problems ranged from $k = 45$ to $k = 85$ by Euclidean norm. The EMD preferred a smaller $k$ ($20$ to $30$), since it is more sensitive to unwanted noise in the reconstruction, while the AUC-IoU generally preferred larger $k$ ($\ge 60$), since it focuses on segmentation overlap, which is less affected by noise.


\subsection{Dynamical low-rank approximation}


\subsubsection{Convergence analysis of DLRA-CG}
\textcolor{red}{TODO:} Not done yet. Some ideas I have:
\begin{itemize}
    \item Connect to CG theory: convergence depends on $\kappa(H)$.
    \item $\Phi$ is quadratic, $H$ is SPD: CG is superior to steepest-descent/Adam (transfers to DLRA-CG, but $\Phi$ isn't quadratic on the manifold). First-order methods don't exploit curvature.
    \item Riemannian GD theory / DLRA theory: difficult, maybe I can get help from Jonas (do we have any convergence guarantees on $\mathcal{M}_r$? Convergence rate?)
    \item The observed gap (50 vs 20'000 iterations) is empirical evidence (I shouldn't over claim theory)
\end{itemize}


\subsubsection{Semi-convergence and restarting}
\textcolor{red}{TODO:} Not done yet. I will discuss the loss of $H$-conjugacy. Some ideas:
\begin{itemize}
    \item Same theory as restarting CG: recover $H$-conjugacy.
    \item If $N_r = 1$ we just have steepest descent.
    \item If $N_r$ is small: smaller "safer" steps (less accumulated conjugacy)
    \item Would be interesting to estimate how fast conjugacy is lost (open question).
    \item Restarting: better recovers true amplitude of source(s), making restarting preferable
\end{itemize}


\subsubsection{Rank as a regularization parameter}
\textcolor{red}{TODO:} Not done yet. We saw that low-rank solutions outperform the exact full-rank weighted Tikhonov. 

Ideas: 
\begin{itemize}
    \item Manifold $\mathcal{M}_r$: forcing $X \in \mathcal{M}_r$ is natural since the true sparse solutions lives on/ near this manifold
    \item If $r$ is chosen too small, it is mathematically impossible to represent some solutions
    \item At rank 5, the solutions resembles the full-rank Tikhonov (a single Gaussian on a 2D mesh is approximately rank-1 as a matrix, a simple outer-product)
    \item Elvetun and Nielsen already did sparsity regularization ($\ell_1$), mention? This method is arguably easier
\end{itemize}


\subsubsection{Preconditioning and local minimas}
We saw precondition sped up convergence, but also drove error up. IC with restarting converged to a "local minima".

\textcolor{red}{TODO:} Not done yet. Ideas:
\begin{itemize}
    \item Classical theory of preconditioning (lower $\kappa$, better convergence)
    \item In the DLRA setting: larger steps on $\mathcal{M}_r$ aer "dangerous". If we get too close to the full-rank solution, we are just making truncations
    \item If we walk directly to (or close to) the full-rank Tikhonov solution, we are in practice just performing a rank $r$ truncation on it
    \item Hard to say why DLRA-PCG using IC and restarting converges to the wrong solution. Maybe just unlucky (this seems to be a local minima: when this solution is used as $X^{(0)}$ for DLRA-CG it also gets stuck here.)
    \item $\mathcal{M}_r$ is non-convex (need reference)
\end{itemize}


%\subsection{Limitations} (Optional)